% GENERATED FILE. Edit canonical lessons, then run npm run generate:books.
% canonical-source-hash: fnv1a64:82df9bbc7431b19a
% canonical-lessons: UR-C244-chhala, UR-C244-nishan, UR-C244-patti, UR-C244-marham, UR-C244-dava-khana

\chapter{Blister, Mark, Bandage, Ointment, Pharmacy}
\label{ch:ur-blister-mark-bandage-ointment-pharmacy}
\hlchaptermodality{244}

\begin{chapteropening}
\textbf{By the end of this chapter:} \emph{I can say a blister, a mark, a bandage, an ointment and a pharmacy.}

Complete the last lesson of chapter 244: I can say a blister, a mark, a bandage, an ointment and a pharmacy.
\end{chapteropening}

\section[\texorpdfstring{\ur{چھالا} (chhālā)}{چھالا (chhālā)}]{\ur{چھالا} (chhālā) — a blister}

\label{lesson:UR-C244-chhala}



\begin{quote}
Before the new one: say the Urdu for height, then the Urdu for a clinic.
\end{quote}

\subsection*{You'll want to know: \ur{چھالا}}
\textbf{\ur{چھالا}} — \emph{chhālā} — \textquotedblleft{}a blister\textquotedblright{}.

\begin{cousinweb}
\textbf{\ur{چھالا}} (\emph{chhālā}) belongs with Hindi \emph{chālā}, \textquotedblleft{}skin, hide\textquotedblright{}, and \emph{chālī}, \textquotedblleft{}a small blister\textquotedblright{}, which CDIAL lists under Sanskrit \emph{challi}, \textquotedblleft{}bark\textquotedblright{}, a word it thinks probably non-Aryan (CDIAL 5005).
\end{cousinweb}

\begin{grammarlens}[title={What you've built}]
One more word for this chapter.
\end{grammarlens}

\subsection*{Your turn}
\begin{itemize}
\raggedright
  \item \emph{Say it:} \emph{chhālā}
  \item \emph{Say it:} \emph{chhālā}, once more
  \item \emph{Say it:} say \emph{chhālā}
  \item \emph{Recall:} say the Urdu for height, then the Urdu for a clinic, then say \emph{chhālā} again
\end{itemize}

\begin{tcolorbox}[breakable,colback=teal!4,colframe=teal!35!black,title={Before you move on}]
What does \ur{چھالا} mean? (\textquotedblleft{}A blister\textquotedblright{}.) Say it once more.
\end{tcolorbox}
\section[\texorpdfstring{\ur{نشان} (nishān)}{نشان (nishān)}]{\ur{نشان} (nishān) — a mark, a scar}

\label{lesson:UR-C244-nishan}



\begin{quote}
Before the new one: say the Urdu for a clinic, then the Urdu for a blister.
\end{quote}

\subsection*{You'll want to know: \ur{نشان}}
\textbf{\ur{نشان}} — \emph{nishān} — \textquotedblleft{}a mark, a scar\textquotedblright{}.

\begin{cousinweb}
\textbf{\ur{نشان}} (\emph{nishān}) is a loan, usually traced to Persian \emph{nishān}, \textquotedblleft{}a sign, a mark\textquotedblright{}.
\end{cousinweb}

\begin{grammarlens}[title={What you've built}]
One more word for this chapter.
\end{grammarlens}

\subsection*{Your turn}
\begin{itemize}
\raggedright
  \item \emph{Say it:} \emph{nishān}
  \item \emph{Say it:} \emph{nishān}, once more
  \item \emph{Say it:} say \emph{nishān}
  \item \emph{Recall:} say the Urdu for a clinic, then the Urdu for a blister, then say \emph{nishān} again
\end{itemize}

\begin{tcolorbox}[breakable,colback=teal!4,colframe=teal!35!black,title={Before you move on}]
What does \ur{نشان} mean? (\textquotedblleft{}A mark, a scar\textquotedblright{}.) Say it once more.
\end{tcolorbox}
\section[\texorpdfstring{\ur{پٹی} (paṭṭī)}{پٹی (paṭṭī)}]{\ur{پٹی} (paṭṭī) — a bandage}

\label{lesson:UR-C244-patti}



\begin{quote}
Before the new one: say the Urdu for a blister, then the Urdu for a mark.
\end{quote}

\subsection*{You'll want to know: \ur{پٹی}}
\textbf{\ur{پٹی}} — \emph{paṭṭī} — \textquotedblleft{}a bandage\textquotedblright{}.

\begin{cousinweb}
\textbf{\ur{پٹی}} (\emph{paṭṭī}) is a tadbhava, a word inherited from Sanskrit and reshaped by sound change. It comes from \emph{paṭṭa}, \textquotedblleft{}cloth; bandage\textquotedblright{} (CDIAL 7700). CDIAL thinks the word probably came from a non-Aryan source.
\end{cousinweb}

\begin{grammarlens}[title={What you've built}]
One more word for this chapter.
\end{grammarlens}

\subsection*{Your turn}
\begin{itemize}
\raggedright
  \item \emph{Say it:} \emph{paṭṭī}
  \item \emph{Say it:} \emph{paṭṭī}, once more
  \item \emph{Say it:} say \emph{paṭṭī}
  \item \emph{Recall:} say the Urdu for a blister, then the Urdu for a mark, then say \emph{paṭṭī} again
\end{itemize}

\begin{tcolorbox}[breakable,colback=teal!4,colframe=teal!35!black,title={Before you move on}]
What does \ur{پٹی} mean? (\textquotedblleft{}A bandage\textquotedblright{}.) Say it once more.
\end{tcolorbox}
\section[\texorpdfstring{\ur{مرہم} (marham)}{مرہم (marham)}]{\ur{مرہم} (marham) — an ointment}

\label{lesson:UR-C244-marham}



\begin{quote}
Before the new one: say the Urdu for a mark, then the Urdu for a bandage.
\end{quote}

\subsection*{You'll want to know: \ur{مرہم}}
\textbf{\ur{مرہم}} — \emph{marham} — \textquotedblleft{}an ointment\textquotedblright{}.

\begin{cousinweb}
\textbf{\ur{مرہم}} (\emph{marham}) is a loan from Arabic \emph{marham}, \textquotedblleft{}ointment\textquotedblright{}, which reached Urdu through Persian. Its Arabic root is \emph{r-h-m} (Lane).
\end{cousinweb}

\begin{grammarlens}[title={What you've built}]
One more word for this chapter.
\end{grammarlens}

\subsection*{Your turn}
\begin{itemize}
\raggedright
  \item \emph{Say it:} \emph{marham}
  \item \emph{Say it:} \emph{marham}, once more
  \item \emph{Say it:} say \emph{marham}
  \item \emph{Recall:} say the Urdu for a mark, then the Urdu for a bandage, then say \emph{marham} again
\end{itemize}

\begin{tcolorbox}[breakable,colback=teal!4,colframe=teal!35!black,title={Before you move on}]
What does \ur{مرہم} mean? (\textquotedblleft{}An ointment\textquotedblright{}.) Say it once more.
\end{tcolorbox}
\section[\texorpdfstring{\ur{دوا} \ur{خانہ} (davā khānā)}{دوا خانہ (davā khānā)}]{\ur{دوا} \ur{خانہ} (davā khānā) — a pharmacy}

\label{lesson:UR-C244-dava-khana}



\begin{quote}
Before the new one: say the Urdu for a bandage, then the Urdu for an ointment.
\end{quote}

\subsection*{You'll want to know: \ur{دوا} \ur{خانہ}}
\textbf{\ur{دوا} \ur{خانہ}} — \emph{davā khānā} — \textquotedblleft{}a pharmacy\textquotedblright{}.

\begin{cousinweb}
\textbf{\ur{دوا} \ur{خانہ}} (\emph{davā khānā}) is \textquotedblleft{}medicine house\textquotedblright{}: \emph{davā}, \textquotedblleft{}medicine\textquotedblright{}, from Arabic \emph{dawāʾ}, \textquotedblleft{}remedy\textquotedblright{}, which reached Urdu through Persian (Lane, root \emph{d-w-y}), and Persian \emph{khāna}, \textquotedblleft{}house\textquotedblright{}.
\end{cousinweb}

\begin{grammarlens}[title={What you've built}]
One more word for this chapter.
\end{grammarlens}

\subsection*{Your turn}
\begin{itemize}
\raggedright
  \item \emph{Say it:} \emph{davā khānā}
  \item \emph{Say it:} \emph{davā khānā}, once more
  \item \emph{Say it:} say \emph{davā khānā}
  \item \emph{Recall:} say the Urdu for a bandage, then the Urdu for an ointment, then say \emph{davā khānā} again
\end{itemize}

\begin{tcolorbox}[breakable,colback=teal!4,colframe=teal!35!black,title={Before you move on}]
What does \ur{دوا} \ur{خانہ} mean? (\textquotedblleft{}A pharmacy\textquotedblright{}.) Say it once more.
\end{tcolorbox}
